% GENERATED FILE. Edit canonical lessons, then run npm run generate:books.
% canonical-source-hash: fnv1a64:397d7a39ccc81ac0
% canonical-lessons: PT-C428-alojamento, PT-C428-pernoitar, PT-C428-dormida, PT-C428-hospedar, PT-C428-porteiro

\chapter{Pernoitar}
\label{ch:pt-pernoitar}
\hlchaptermodality{428}

\begin{chapteropening}
\textbf{By the end of this chapter:} \emph{I can talk about where I stay: accommodation, spending the night, an overnight stay, putting someone up and the doorman.}

Complete the last lesson of chapter 428: I can talk about where I stay: accommodation, spending the night, an overnight stay, putting someone up and the doorman.
\end{chapteropening}

\section[o alojamento]{o alojamento — the accommodation, the lodging}

\label{lesson:PT-C428-alojamento}



\begin{quote}
Before the new one: say the Portuguese for to bring forward, then the Portuguese for to rebook.
\end{quote}

\subsection*{You'll want to know: o alojamento}
\textbf{o alojamento} — \textquotedblleft{}the accommodation, the lodging\textquotedblright{}.

\begin{quote}
O alojamento perto da praia é caro.
\end{quote}

\begin{quote}
\emph{Accommodation near the beach is expensive.}
\end{quote}

\emph{O alojamento local} — short holiday lets in Portugal.

\begin{cousinweb}
\emph{alojamento} is built on \emph{alojar} 'to lodge', usually traced to French \emph{loger}, on \emph{loge} 'hut, shelter', which EtymDB traces to Frankish \emph{laubija} 'shelter, roof'.
\end{cousinweb}

\begin{grammarlens}[title={What you've built}]
One more word for telling, explaining and describing at B1.
\end{grammarlens}

\subsection*{Your turn}
\begin{itemize}
\raggedright
  \item \emph{Say it:} \emph{o alojamento}
  \item \emph{Say it:} \emph{o alojamento}, once more
  \item \emph{Say it:} a sentence of your own with \emph{o alojamento}
  \item \emph{Recall:} say the Portuguese for to bring forward, then the Portuguese for to rebook, then say \emph{o alojamento} again
\end{itemize}

\begin{tcolorbox}[breakable,colback=teal!4,colframe=teal!35!black,title={Before you move on}]
What does o alojamento mean? (\textquotedblleft{}The accommodation, the lodging\textquotedblright{}.) Say it once more.
\end{tcolorbox}
\section[pernoitar]{pernoitar — to spend the night, to stay overnight}

\label{lesson:PT-C428-pernoitar}



\begin{quote}
Before the new one: say the Portuguese for to rebook, then the Portuguese for the accommodation.
\end{quote}

\subsection*{You'll want to know: pernoitar}
\textbf{pernoitar} — \textquotedblleft{}to spend the night, to stay overnight\textquotedblright{}.

\begin{quote}
Pernoitámos numa aldeia a meio do caminho.
\end{quote}

\begin{quote}
\emph{We spent the night in a village halfway there.}
\end{quote}

A little formal; people also say \emph{passar a noite}.

\begin{cousinweb}
\emph{pernoitar} comes from Latin \emph{pernoctāre} 'to spend the night' (EtymDB), reshaped on \emph{noite} 'night', from Latin \emph{noctem}.
\end{cousinweb}

\begin{grammarlens}[title={What you've built}]
One more word for telling, explaining and describing at B1.
\end{grammarlens}

\subsection*{Your turn}
\begin{itemize}
\raggedright
  \item \emph{Say it:} \emph{pernoitar}
  \item \emph{Say it:} \emph{pernoitar}, once more
  \item \emph{Say it:} a sentence of your own with \emph{pernoitar}
  \item \emph{Recall:} say the Portuguese for to rebook, then the Portuguese for the accommodation, then say \emph{pernoitar} again
\end{itemize}

\begin{tcolorbox}[breakable,colback=teal!4,colframe=teal!35!black,title={Before you move on}]
What does pernoitar mean? (\textquotedblleft{}To spend the night, to stay overnight\textquotedblright{}.) Say it once more.
\end{tcolorbox}
\section[a dormida]{a dormida — the overnight stay; a place to sleep}

\label{lesson:PT-C428-dormida}



\begin{quote}
Before the new one: say the Portuguese for the accommodation, then the Portuguese for to spend the night.
\end{quote}

\subsection*{You'll want to know: a dormida}
\textbf{a dormida} — \textquotedblleft{}the overnight stay; a place to sleep\textquotedblright{}.

\begin{quote}
A dormida custa trinta euros com pequeno-almoço.
\end{quote}

\begin{quote}
\emph{A night costs thirty euros with breakfast.}
\end{quote}

Built on \emph{dormir}.

\begin{cousinweb}
\emph{dormida} is built on \emph{dormir} 'to sleep', which is inherited from Latin \emph{dormīre} (Lewis \& Short).
\end{cousinweb}

\begin{grammarlens}[title={What you've built}]
One more word for telling, explaining and describing at B1.
\end{grammarlens}

\subsection*{Your turn}
\begin{itemize}
\raggedright
  \item \emph{Say it:} \emph{a dormida}
  \item \emph{Say it:} \emph{a dormida}, once more
  \item \emph{Say it:} a sentence of your own with \emph{a dormida}
  \item \emph{Recall:} say the Portuguese for the accommodation, then the Portuguese for to spend the night, then say \emph{a dormida} again
\end{itemize}

\begin{tcolorbox}[breakable,colback=teal!4,colframe=teal!35!black,title={Before you move on}]
What does a dormida mean? (\textquotedblleft{}The overnight stay; a place to sleep\textquotedblright{}.) Say it once more.
\end{tcolorbox}
\section[hospedar]{hospedar — to put up, to host (guests)}

\label{lesson:PT-C428-hospedar}



\begin{quote}
Before the new one: say the Portuguese for to spend the night, then the Portuguese for the overnight stay.
\end{quote}

\subsection*{You'll want to know: hospedar}
\textbf{hospedar} — \textquotedblleft{}to put up, to host (guests)\textquotedblright{}.

\begin{quote}
Os meus tios hospedaram-nos durante uma semana.
\end{quote}

\begin{quote}
\emph{My aunt and uncle put us up for a week.}
\end{quote}

\emph{Hospedar-se} — to stay: \emph{hospedámo-nos num hotel}.

\begin{cousinweb}
\emph{hospedar} is inherited from Late Latin \emph{hospitāre} 'to receive as a guest' (EtymDB), built on \emph{hospes} 'host, guest'.
\end{cousinweb}

\begin{grammarlens}[title={What you've built}]
One more word for telling, explaining and describing at B1.
\end{grammarlens}

\subsection*{Your turn}
\begin{itemize}
\raggedright
  \item \emph{Say it:} \emph{hospedar}
  \item \emph{Say it:} \emph{hospedar}, once more
  \item \emph{Say it:} a sentence of your own with \emph{hospedar}
  \item \emph{Recall:} say the Portuguese for to spend the night, then the Portuguese for the overnight stay, then say \emph{hospedar} again
\end{itemize}

\begin{tcolorbox}[breakable,colback=teal!4,colframe=teal!35!black,title={Before you move on}]
What does hospedar mean? (\textquotedblleft{}To put up, to host (guests)\textquotedblright{}.) Say it once more.
\end{tcolorbox}
\section[o porteiro]{o porteiro — the doorman, the porter}

\label{lesson:PT-C428-porteiro}



\begin{quote}
Before the new one: say the Portuguese for the overnight stay, then the Portuguese for to put up.
\end{quote}

\subsection*{You'll want to know: o porteiro}
\textbf{o porteiro} — \textquotedblleft{}the doorman, the porter\textquotedblright{}.

\begin{quote}
O porteiro do hotel chamou um táxi.
\end{quote}

\begin{quote}
\emph{The hotel doorman called a taxi.}
\end{quote}

In a block of flats, \emph{o porteiro} looks after the entrance.

\begin{cousinweb}
\emph{porteiro} is inherited (EtymDB) from Latin \emph{portārius} 'doorkeeper', built on \emph{porta} 'door' (Lewis \& Short).
\end{cousinweb}

\begin{grammarlens}[title={What you've built}]
One more word for telling, explaining and describing at B1.
\end{grammarlens}

\subsection*{Your turn}
\begin{itemize}
\raggedright
  \item \emph{Say it:} \emph{o porteiro}
  \item \emph{Say it:} \emph{o porteiro}, once more
  \item \emph{Say it:} a sentence of your own with \emph{o porteiro}
  \item \emph{Recall:} say the Portuguese for the overnight stay, then the Portuguese for to put up, then say \emph{o porteiro} again
\end{itemize}

\begin{tcolorbox}[breakable,colback=teal!4,colframe=teal!35!black,title={Before you move on}]
What does o porteiro mean? (\textquotedblleft{}The doorman, the porter\textquotedblright{}.) Say it once more.
\end{tcolorbox}
